% ----------------------
\subsection{Section 32 (October 1830)}
% ----------------------
	\label{DC32}
	
	% -----
	\subsubsection{JSP version of this section}
	% -----
	
		From the JSP: ``\textit{Revelation, Manchester Township, Ontario Co., NY, Oct. 1830. Featured version copied [ca. Aug. 1834] in Revelation Book 2, pp. 83--84; handwriting of Frederick G. Williams; Revelations Collection, CHL.}
		
		``For unknown reasons, this revelation was not copied into Revelation Book 1, nor was it published in the 1833 Book of Commandments. Therefore, this version, created mid-1834, is the earliest extant version of this revelation.''
		
		See below in the historical introduction for more discussion.

	% -----
	\subsubsection{Historical Background}
	% -----
		\label{DC32-HB}
		
		% --
		\paragraph{JSP}
		% --
		
			``After Oliver Cowdery and Peter Whitmer Jr. were called on a mission to preach to the American Indians in September 1830, several other elders expressed a desire to see the fulfillment of God's purposes among `the remnants of the house of Joseph---the Lamanites residing in the west.' JS inquired `of the Lord respecting the propriety of sending som[e] of the Elders among them,' resulting in the following revelation for Parley P. Pratt and Ziba Peterson.

			``With his longstanding interest in the American Indians, Pratt no doubt welcomed the call. He had set out for Ohio from New York in the fall of 1826, when he was nineteen, with plans to preach to the Indians. He later recalled his youthful thoughts: `I will win the confidence of the red man; I will learn his language; I will tell him of Jesus; I will read to him the Scriptures; I will teach him the arts of peace; to hate war, to love his neighbor, to fear and love God.'3 However, Pratt did not then follow through with his plan. Instead, he built a small log hut some thirty miles west of Cleveland, Ohio, and spent the winter in reading and study, especially of the Bible. The following summer he returned to New York, married Thankful Halsey, and went back to Ohio with his new bride. In 1829, Pratt, who was a Baptist, was influenced by Sidney Rigdon, a restorationist preacher affiliated with Alexander Campbell. Attracted by the Campbellites' efforts to recreate the primitive church of the New Testament and by their expectation of Christ's second coming, Pratt decided to become an itinerant preacher of their message. In the summer of 1830, Pratt and his wife left Ohio for their `native place,' Columbia County, New York, where he planned to preach. En route by canal boat, Pratt responded to what he felt were spiritual promptings and disembarked at Newark, near Palmyra, leaving his wife to travel on alone. Pratt encountered a Baptist deacon who loaned him a copy of the Book of Mormon; he soon `knew and comprehended that the book was true.' Pratt then traveled to the Palmyra area, where he met Hyrum Smith. Shortly after, the two traveled to Fayette, where Pratt was baptized into the Church of Christ and ordained an elder. He then continued on to Columbia County to rejoin his wife and proselytize friends and family.
			
			``Pratt soon learned of the call Cowdery and Whitmer had received to preach to the Lamanites. Not long after, the revelation featured here called Pratt, along with Ziba Peterson, to accompany them. Peterson, a resident of Macedon, New York, was a founding elder of the Church of Christ and received an elder's license on 9 June 1830.
			
			``Unlike other revelations from this period, this revelation was not recorded in Revelation Book 1. It is possible that Pratt took the only copy with him when he left on his mission, making it unavailable for inclusion when John Whitmer was compiling the relevant portion of the book. When this revelation was later recorded in Revelation Book 2 (the source for the text presented here), it included the notation `Manchester Oct 1830.' Though the precise date of the revelation was not recorded, it was likely dictated shortly before 17 October, when Pratt and the other missionaries assigned to preach to the Lamanites signed a covenant to follow Cowdery as he led them in their missionary labors.''
			
	% -----
	\subsubsection{Versions}
	% -----
		\label{DC32-V}
		
		See the \href{https://drive.google.com/file/d/1goAvNz8jS4R8slYfMG_db55Ty2z_vmgK/view?usp=sharing}{sections comparison} document.
	
		\begin{compactdesc}
			\item[JSP] \hyperref[EMS-1830-JS-OCT-1830-A]{October 1830}
			\item[BC] ---
			\item[DC 1835] Section 54, 184
			\item[TS] 4:11 (15 April 1843), 172
			\item[MS] 4:11 (March 1844), 168
			\item[DC 1844] Section 54, 277
		\end{compactdesc}

	% -----
	\subsubsection{Section 32:1} % *DC32-1 
	% -----
		\label{DC32-1}

		AND now concerning my servant Parley P. Pratt, behold, I say unto him that as I live I will that he shall declare my gospel and learn of me, and be meek and lowly%
			\footnote{%
				% \label{}%
				For other occurrences of ``meek and lowly,'' see \hyperref[MEEKNESS-PHRASES]{here}.
				}
		of heart.

		\begin{framed}
		\scriptsize
			\begin{compactdesc}
				\item[JSP] And now concerning my servant Parley behold I say unto him that as I live I will that he shall declare my gospel and Learn of me and be meek and lowly of heart
				\item[BC] ---
				% \item[]
			\end{compactdesc}
		\end{framed}

	% -----
	\subsubsection{Section 32:2} % *DC32-2 
	% -----
		\label{DC32-2}

		And that which I have appointed unto him is that he shall go with my servants, Oliver Cowdery and Peter Whitmer, Jun., into the wilderness among the Lamanites.

		\begin{framed}
		\scriptsize
			\begin{compactdesc}
				\item[JSP] and that which I have appointed unto him is that he shall go with my servant Oliver [Cowdery] and Peter [Whitmer Jr.] into the wilderness among the Lamanites
				\item[BC] ---
				% \item[]
			\end{compactdesc}
		\end{framed}

	% -----
	\subsubsection{Section 32:3} % *DC32-3 
	% -----
		\label{DC32-3}

		And Ziba Peterson also shall go with them; and I myself will go with them and be in their midst; and I am their \hyperref[ADVOCATE]{advocate} with the Father, and nothing shall prevail against them.

		\begin{framed}
		\scriptsize
			\begin{compactdesc}
				\item[JSP] and Ziba [Peterson] also shall go with them and I myself will go with them and be in their midst and I am their advocate with the Father and nothing shall prevail
				\item[BC] ---
				% \item[]
			\end{compactdesc}
		\end{framed}

	% -----
	\subsubsection{Section 32:4} % *DC32-4 
	% -----
		\label{DC32-4}

		And they shall give heed to that which is written, and pretend to%
			\footnote{%
				% \label{}%
				1828 definition \#5 for ``pretend'' is ``to claim,'' with the note that ``in this we generally use \textit{pretend to}.''
				}
		no other revelation; and they shall pray always that I may \hyperref[UNFOLD]{unfold} the same to their understanding.

		\begin{framed}
		\scriptsize
			\begin{compactdesc}
				\item[JSP] and they shall give heed to that which is writen and pretend to no other revelation and they shall pray always that I may unfold them to their understanding
				\item[BC] ---
				% \item[]
			\end{compactdesc}
		\end{framed}

	% -----
	\subsubsection{Section 32:5} % *DC32-5 
	% -----
		\label{DC32-5}

		And they shall give heed unto these words and trifle not,%
			\footnote{%
				% \label{}%
				``Trifle'' appears four times in the scriptures: here, \hyperref[MOSIAH-2-9]{Mosiah 2:9}, \hyperref[DC6-12]{DC 6:12}, and \hyperref[DC8-10]{DC 8:10}. From the 1828 dictionary for ``trifle'' as intransitive verb:
					\begin{compactitem}
						\item To act or talk without seriousness, gravity, weight or dignity; to act or talk with levity.
						\item[1] To indulge in light amusements.
						\item To \textit{trifle} with, to mock; to play the fool with; to treat without respect or seriousness.
						\item To \textit{trifle} with, to spend in vanity; to waste.
						\item To \textit{trifle} away, to no good purpose; as, to \textit{trifle} with time, or to \textit{trifle} away time; to \textit{trifle} with advantages.
					\end{compactitem}
				}
		and I will bless them.  Amen.

		\begin{framed}
		\scriptsize
			\begin{compactdesc}
				\item[JSP] and they shall give heed unto these words and trifle not and I will bless them amen Manchester Oct 1830---
				\item[BC] ---
				% \item[]
			\end{compactdesc}
		\end{framed}